\section{Data and methods}\label{sec:data-methods}

\subsection{Observations}
% P:methods-01
All solutions are 30 wt\% \MEA on a \COtwo-free basis unless stated otherwise.
For the Wagner comparisons, each row is evaluated at its reported temperature and \COtwo-free \MEA mass fraction: 29.93--30.52 wt\% near \SI{353}{\kelvin} and 30.28--30.95 wt\% near \SI{392}{\kelvin}.
Supplementary Section~S2 places the observations below within an inventory of published measurements near 30 wt\% \MEA.
The pressure comparisons use 161 measured \COtwo partial pressures at 40--\SI{120}{\degreeCelsius} from six studies \cite{Aronu2011,Hilliard2008,idrisSpeciationMEACO2Adducts2014,Jou1995,Mamun2005,Xu2011}.
Each row keeps its source table and row, the reported temperature, the loading and the pressure in kilopascal.
Six Xu et al.\ rows, tabulated at measured temperatures of 100.5--\SI{101.1}{\degreeCelsius} and 120.4--\SI{121.8}{\degreeCelsius}, are grouped with and calculated at the 100 and \SI{120}{\degreeCelsius} isotherms in this work, following the comparison temperatures of the source's Fig.~3 \cite[Table~1, Fig.~3]{Xu2011}.
Idris et al.\ report a combined standard uncertainty for each pressure \cite[Table~2]{idrisSpeciationMEACO2Adducts2014}, and the other five studies give no row-level uncertainty that we retained.

% P:methods-02
The liquid species data come from two methods.
B\"ottinger et al.\ measured species by online \textsuperscript{1}H and \textsuperscript{13}C NMR spectroscopy \cite[Sect.~3]{Bottinger2008}.
Fast proton exchange makes only the sum \ce{MEA + MEAH+} observable, and the authors report carbonate as practically absent, so their bicarbonate value is compared with the modeled \ce{HCO3- + CO3^2-} pool.
Matin et al.\ derived species from acid titration for total alkalinity, strong-base titration and total inorganic carbon \cite{matinFacileMethodDetermination2012}.
Their method reports \MEA, \MEAH and \MEACOO separately and assigns carbonate to bicarbonate, so their bicarbonate value is also compared with the \ce{HCO3- + CO3^2-} pool \cite[Eqs.~18--19b]{matinFacileMethodDetermination2012}.
Matin et al.\ label their measurements \SI{21}{\degreeCelsius}, and we represent them at \SI{20}{\degreeCelsius}.
Jakobsen et al.\ report bicarbonate and carbonate separately \cite{Jakobsen2005}; their observations enter the additional species comparisons in \cref{fig:speciation} and the carbonate-share comparison in the Results, but not the present fit.
Every species target is a true-species liquid mole fraction.

% P:methods-03
\Cref{tab:mea-data-inventory} summarizes the fitted target sets and comparison groups; these groups overlap and their counts are not additive.
The fitted and 80\,\si{\degreeCelsius} targets belong to an estimation set of 123 states: 79 pressure states from Hilliard and Jou et al.\ and 44 species states from B\"ottinger and Matin et al.
Three species states are excluded from the objective; their observations are compared but not fitted.
Two kinds of fixed values described below were estimated partly on 80\,\si{\degreeCelsius} data.
The 80\,\si{\degreeCelsius} targets therefore test conditions outside the present fit, not data unseen during model development.

\begin{table*}[t]
    \centering
    \caption{Observation roles.
    The pool is the \ce{HCO3- + CO3^2-} mole fraction.
    Fitted and 80\,\si{\degreeCelsius} targets belong to the 123-state estimation set; the pooled and out-of-range rows belong to the 161 measured pressures at 40--\SI{120}{\degreeCelsius}.}
    \label{tab:mea-data-inventory}
    \begin{tabularx}{\textwidth}{@{}>{\raggedright\arraybackslash}p{3.3cm}r>{\raggedright\arraybackslash}X>{\raggedright\arraybackslash}p{4.6cm}@{}}
        \toprule
        Role & Targets & Sources and conditions & Observable \\
        \midrule
        Fitted pressure & 48 & Hilliard 31, Jou et al.\ 17; 32 at 40 and 16 at \SI{60}{\degreeCelsius} & \(p_{\COtwo}\) \\
        Fitted species & 94 & B\"ottinger et al.\ 40 at 20, 40 and \SI{60}{\degreeCelsius}; Matin et al.\ 54 at \SI{20}{\degreeCelsius} & B\"ottinger: \ce{MEA + MEAH+} 18, \MEACOO 18, pool 4; Matin: \MEA, \MEAH and \MEACOO, 18 each \\
        Species compared, not fitted & 18 & Matin et al., \SI{20}{\degreeCelsius} & Pool \\
        Out-of-fit, \SI{80}{\degreeCelsius} & 11 + 11 & Jou et al.\ pressure; B\"ottinger et al.\ species & \(p_{\COtwo}\); \ce{MEA + MEAH+} 5, \MEACOO 5, pool 1 \\
        Pooled pressure & 104 & Measured pressures at 40--\SI{80}{\degreeCelsius}: Aronu 36, Hilliard 30, Jou 28, Idris 10 & \(p_{\COtwo}\) \\
        Previously unused pressure & 11 + 12 & Wagner et al.\ near 353 and \SI{392}{\kelvin}, used in no fit or selection \cite{wagnerSolubilityCarbonDioxide2013} & \(p_{\COtwo}\) \\
        Out-of-range pressure & 57 & Measured pressures at 100--\SI{120}{\degreeCelsius}: Jou, Mamun and Xu et al. & \(p_{\COtwo}\) \\
        Concentration transfer & 70 & Aronu et al.\ at 15 wt\% (33) and 45 wt\% (37), 40--\SI{80}{\degreeCelsius} & \(p_{\COtwo}\) \\
        Carbonate comparison & 10 & Jakobsen et al.\ & Carbonate fraction of dissolved carbon \\
        \bottomrule
    \end{tabularx}
\end{table*}

\subsection{Parameters}
% P:methods-04
\Cref{tab:full-ionic-ssm-ds-parameters,tab:binary-interaction-parameters,tab:reaction-equilibrium-constants} list every parameter of the fitted SSM+DS parameter set, called the fitted parameter set below.
Only the five binary interaction parameters of \cref{tab:binary-interaction-parameters} are fitted here, to the 142 targets of \cref{tab:mea-data-inventory}.
The molecular and \HthreeO, \Hydroxide, \HCOthree and \COthree parameters come from the cited sources, with the exceptions marked in the tables.
The following quantities were estimated from data while the model was developed and are held fixed here; they are fitted or data-selected values, not independently estimated molecular properties.
The segment diameters, dispersion energies and Born diameters of \MEAH and \MEACOO were fitted to liquid speciation alone: NMR values of B\"ottinger et al.\ and Jakobsen et al.\ and titration values of Matin et al.\ at 20--\SI{60}{\degreeCelsius}~\cite{Bottinger2008,Jakobsen2005,matinFacileMethodDetermination2012}.
The \HCOthree and \COthree Born diameters were kept at \SI{3.0}{\angstrom}: a fit of both diameters to the bicarbonate and carbonate observations of Jakobsen et al.\ and Matin et al.\ at 20 and \SI{40}{\degreeCelsius} stayed at that starting value, and starts at other values ended at different diameters without converging~\cite{Jakobsen2005,matinFacileMethodDetermination2012}.
Their selection history is only partly retained, and they are neither independently determined ion properties nor uniquely identified fitted diameters.
The \COtwo dispersion energy was selected by comparison against 161 pressure measurements and 131 species observations that included the 80\,\si{\degreeCelsius} isotherm.
The R5 coefficient \(A_5\) and the R2 and R5 reaction enthalpies were changed as described below, using data that included the 80\,\si{\degreeCelsius} isotherm.
The \MEA--water and \COtwo--water interactions were fitted to binary VLE at 363--\SI{432}{\kelvin} and 313--\SI{393}{\kelvin} \cite{caiBinaryIsobaricVaporLiquid1996,kiepeExperimentalDeterminationPrediction2002}.

% P:methods-05
The R2 and R5 coefficients carry fixed constant shifts of the reaction enthalpy, \SI{-8.2019}{\kilo\joule\per\mole} for R2 and \SI{8.4185}{\kilo\joule\per\mole} for R5, applied so that \(\ln K\) is unchanged at \SI{313.15}{\kelvin}.
They were estimated against pressure, species and heat-of-absorption targets that included the 80\,\si{\degreeCelsius} isotherm, among them two of the 11 Jou et al.\ measurements at \SI{80}{\degreeCelsius}, at loadings 0.118 and 0.578.
The R5 shift was applied to an \(A_5\) of \SI{2597.91}{\kelvin}, \SI{80}{\kelvin} below the Bates--Pinching value of \SI{2677.91}{\kelvin}.
This lower \(A_5\) lowers \(\mathrm{p}K_5=-\log_{10}K_5\) by \((\SI{80}{\kelvin})/T\) at every temperature, so \MEAH dissociates more readily; it was chosen by comparing fixed-parameter SSM+DS calculations with pressure and species data that included the 80\,\si{\degreeCelsius} isotherm.
The same estimation also shifted R4, but the fitted parameter set uses the R4 source correlation.

\begin{table*}[t]
    \centering
    \caption{Component parameters of the selected SSM+DS record.
    All values are held fixed in the 142-target fit.
    Ion hard-chain and Debye--H\"uckel diameters are \(0.88\,\sigma_i\).
    The water segment diameter is \(\sigma_{\HtwoO}/\text{\AA}=2.7927+10.11\exp(-0.01775\,T/\mathrm{K})-1.417\exp(-0.01146\,T/\mathrm{K})\) \cite{Held2008}, and its permittivity is \cref{eq:water-permittivity}.
    Association columns give the self-association volume \(\kappa^{AB}\) (\(\sigma^3\) convention) and energy \(\epsilon^{AB}/k_{\mathrm{B}}\) of the 2B scheme; \COtwo has two induced-association sites that bond only with water (\cref{sec:theory}).
    A dash marks a parameter that the component does not have.}
    \label{tab:full-ionic-ssm-ds-parameters}
    \footnotesize\rmfamily
    \setlength{\tabcolsep}{3pt}
    \renewcommand{\arraystretch}{1.12}
    \begin{tabularx}{\textwidth}{@{}lrrrrrrrrrr>{\raggedright\arraybackslash}X@{}}
        \toprule
        Species & \(z_i\) & \makecell{\(M_i\)\\(kg/mol)} & \(m_i\) & \makecell{\(\sigma_i\)\\(\AA)} & \makecell{\(u_i/k_{\mathrm{B}}\)\\(K)} & \makecell{\(d_i^{\mathrm{Born}}\)\\(\AA)} & \(f_k\) & \(\varepsilon_i\) & \(\kappa^{AB}\) & \makecell{\(\epsilon^{AB}/k_{\mathrm{B}}\)\\(K)} & Source and role \\
        \midrule
        \COtwo & 0 & 0.04401 & 2.0729 & 2.7852 & 173.440 & -- & 1 & -- & -- & -- & \(m\), \(\sigma\): \cite{Gross2001}; \(u/k_{\mathrm{B}}\) inherited from a 2026 calibration, 2.5\,\% above 169.21 \cite{Gross2001}; induced association \cite{Schick2023} \\
        \MEA & 0 & 0.06108 & 3.0353 & 3.0435 & 277.174 & -- & 1 & 32 & 0.036287 & 2586.3 & \cite[Table~2]{Baygi2015}; \(\kappa^{AB}\) converted from 0.037470 \\
        \HtwoO & 0 & 0.01801528 & 1.2047 & \(\sigma_{\HtwoO}(T)\) & 353.95 & -- & 1.5 & \(\varepsilon_{\HtwoO}(T)\) & 0.04509 & 2425.7 & \cite{Held2008,Cameretti2005}; \(f_{\HtwoO}\) \cite[Table~2]{Figiel2025} \\
        \MEAH & \(+1\) & 0.06209 & 1 & 3.4851 & 232.687 & 3.5332 & 1 & -- & -- & -- & Inherited from an earlier project speciation fit \\
        \MEACOO & \(-1\) & 0.10408 & 1 & 3.5354 & 453.265 & 3.5411 & 1 & -- & -- & -- & Inherited from an earlier project speciation fit \\
        \HCOthree & \(-1\) & 0.0610168 & 1 & 2.9296 & 70.00 & 3.0 & 1 & -- & -- & -- & \(\sigma\), \(u/k_{\mathrm{B}}\): \cite[Table~2]{Held2014}; \(d^{\mathrm{Born}}\) inherited project value \\
        \COthree & \(-2\) & 0.06001 & 1 & 2.4422 & 249.26 & 3.0 & 1 & -- & -- & -- & \(\sigma\), \(u/k_{\mathrm{B}}\): \cite[Table~2]{Held2014}; \(d^{\mathrm{Born}}\) inherited project value \\
        \HthreeO & \(+1\) & 0.01902 & 1 & 3.4654 & 500.00 & 1.218 & 1 & -- & -- & -- & \cite[Table~2]{Held2014}; \(d^{\mathrm{Born}}\) \cite[Table~3]{Figiel2025} \\
        \Hydroxide & \(-1\) & 0.01701 & 1 & 2.0177 & 650.00 & 3.0811 & 1 & -- & -- & -- & \(\sigma\), \(u/k_{\mathrm{B}}\): \cite[Table~2]{Held2014}; \(d^{\mathrm{Born}}\) from a Born hydration-energy inversion \\
        \bottomrule
    \end{tabularx}
\end{table*}

% Generated by scripts/final_rerun_tables.py from analyses/mea_parameter_bundle/results/heat-final-170; do not hand-edit values.
\begin{table*}[t]
    \centering
    \caption{Binary dispersion interaction parameters of the selected SSM+DS fit (F1) and the original Born fit (F2), with conditional standard errors (SE).
    The \MEAH--water coefficient is the value at \(T_{\mathrm{ref}}=\SI{313.15}{\kelvin}\) in \cref{eq:kij-temperature}, and \(b\) is its reciprocal-temperature slope.
    SEs condition on each fit's active bounds, use the declared residual scales and treat residuals as independent, although the titration species of one sample are not; they are indicative only.
    A coordinate at a bound has no symmetric interval.
    Pairs of ions with the same charge sign have no dispersion interaction, and every pair not listed has \(k_{ij}=0\).
    Full-precision values are in the parameter file named in the Data and Code Availability statement.}
    \label{tab:binary-interaction-parameters}
    \begin{tabularx}{\textwidth}{@{}lllrrrr>{\raggedright\arraybackslash}X@{}}
        \toprule
        Parameter & Role & Bounds & F1 & SE & F2 & SE & Source or note \\
        \midrule
        \(k_{\MEACOO,\HtwoO}\) & Fitted & \(\pm0.5\) & 0.2240 & 0.0118 & 0.0918 & 0.0151 & This work \\
        \(k_{\MEAH,\HtwoO}\) & Fitted & \(\pm0.5\) & \(-\)0.3125 & 0.0128 & 0.1547 & 0.0149 & This work; value at \(T_{\mathrm{ref}}\) \\
        \(b\) (K) & Fitted & \(\pm1000\) & 107.96 & 22.93 & 106.99 & 30.72 & This work \\
        \(k_{\HCOthree,\HtwoO}\) & Fitted & \(\pm0.5\) & \(+\)0.5 (bound) & -- & \(+\)0.5 (bound) & -- & This work \\
        \(k_{\MEAH,\MEACOO}\) & Fitted & \(\pm1\) & \(-\)0.9683 & 0.0267 & \(-\)1 (bound) & -- & This work \\
        \(k_{\MEA,\HtwoO}\) & Fixed & -- & \(-0.0735\) & -- & \(-0.0735\) & -- & Fitted to binary VLE \cite{caiBinaryIsobaricVaporLiquid1996} \\
        \(k_{\COtwo,\HtwoO}\) & Fixed & -- & 0.0133 & -- & 0.0133 & -- & Fitted to binary VLE \cite{kiepeExperimentalDeterminationPrediction2002} \\
        \(k_{\HthreeO,\HtwoO}\) & Fixed & -- & 0.25 & -- & 0.25 & -- & \cite[Table~2]{Held2014} \\
        \(k_{\Hydroxide,\HtwoO}\) & Fixed & -- & \(-0.25\) & -- & \(-0.25\) & -- & \cite[Table~2]{Held2014} \\
        \(k_{\COthree,\HtwoO}\) & Fixed & -- & \(-0.25\) & -- & \(-0.25\) & -- & \cite[Table~2]{Held2014} \\
        \bottomrule
    \end{tabularx}
\end{table*}


\subsection{Objective and estimation}
% P:methods-06
The fit minimizes \(C=\tfrac12\sum_jr_j^2\) over the 142 fitted targets, with
\begin{equation}
    r_{P,j}=\frac{\ln\left(p_{\COtwo,j}^{\mathrm{calc}}/p_{\COtwo,j}^{\mathrm{obs}}\right)}{0.3},
    \qquad
    r_{x,j}=\frac{x_j^{\mathrm{calc}}-x_j^{\mathrm{obs}}}{0.1\,x_j^{\mathrm{obs}}+0.001}.
    \label{eq:fit-residuals}
\end{equation}
These scales are the declared fitting weights, not measured standard deviations, and every target has unit weight within its family.
The bounds are \(\pm0.5\) for the three ion--water coefficients, \(\pm1\) for \MEAH--\MEACOO and \(\pm\SI{1000}{\kelvin}\) for \(b\), with parameter scales of 0.01 and \SI{10}{\kelvin}.
The same binary interaction parameters, bounds, weights, reactions and ion inputs define the SSM+DS and original Born fits.

% P:methods-07
Two choices of the objective were made with knowledge of the fits that preceded them.
The 100--\SI{120}{\degreeCelsius} pressure targets are excluded from the objective.
The 18 Matin pool targets are compared but not fitted, on the basis of the authors' own comparison with NMR data \cite{matinFacileMethodDetermination2012}.
Two sensitivity fits use the same parameters, bounds and weights: the fit including the titration bicarbonate data adds these 18 targets to the objective, and the fit without the titration data omits all 72 Matin targets.

% P:methods-08
Each fit is a bounded Levenberg--Marquardt least-squares problem with exact implicit derivatives of the equilibrium states with respect to the fitted parameters.
The SSM+DS fit including the titration bicarbonate data used six prescribed starts; five converged, and one failed its initial equilibrium evaluation.
The fit to the 142 targets started from the best of those optima and stopped when the relative cost change between updates fell below the \(10^{-12}\) tolerance; the final change was \(2.6\times10^{-13}\), with no failed trial state.
This is a local optimum, not evidence of a global one.
Local identification uses the column-normalized singular values of the weighted Jacobian at the solution.
No covariance is reported because the \HCOthree--water coefficient lies on its bound.

\subsection{Numerical verification and metrics}
% P:methods-09
The software builds used for each calculation are identified in the Data and Code Availability statement.
In every equilibrium calculation compared with pressure or species observations, the \COtwo-equivalent carbon of the feed equals the reported loading, \(\alpha=(n_{\mathrm{C}}-2n_{\mathrm{am}})/n_{\mathrm{am}}\), where \(n_{\mathrm{C}}\) is the total carbon of the feed, including the small ion amounts that start the solution, and \(n_{\mathrm{am}}\) is the total amine, each \MEA unit carrying two carbons.
A state is accepted when the solver meets its requested tolerance, the material balances close within \(10^{-7}\max(1,|n|)\)~\si{\mole} for each conserved total \(n\), and the charge balance closes within \(10^{-7}\)~mol of elementary charge.
All 123 estimation states and all 161 measured-pressure states pass, with largest absolute stationarity residuals of \(4.0\times10^{-13}\) and \(1.7\times10^{-13}\).
Re-evaluation of the fitted parameter set reproduces the fitted cost within \(4.2\times10^{-12}\).
These checks verify the numerical solutions, not the physical model.

% P:methods-10
Agreement is reported as
\begin{equation}
    \mathrm{AARD}=\frac{100\,\%}{n}\sum_{j=1}^{n}\left|\frac{y_j^{\mathrm{calc}}}{y_j^{\mathrm{obs}}}-1\right|,
    \label{eq:aard}
\end{equation}
with the mean and root mean square of \(\ln(y^{\mathrm{calc}}/y^{\mathrm{obs}})\) reported as signed logarithmic bias and overall logarithmic error, respectively.
Each value is computed from the retained calculations and rounded for display.
